\section{Esito della valutazione}

Il confronto tra i dieci capitolati, integrato dagli incontri con i proponenti, ha portato il gruppo \emph{Git Happens} a candidarsi per il capitolato \textbf{C8 Sonar} di Halue S.r.l. Società Benefit. Esaurita la disponibilità di C0 AIRA, prima preferenza del gruppo, C8 si è rivelato la scelta più equilibrata: a differenza delle proposte vincolate da componenti hardware (C1, C3), da forti vincoli normativi e da un perimetro ancora da definire (C5) o da limiti di dominio, di perimetro o di natura del lavoro (C2, C4, C6, C7, C9), Sonar unisce un problema reale, requisiti chiari e verificabili, uno stack moderno e un proponente disponibile e flessibile, offrendo le condizioni ideali per applicare con successo le metodologie di Ingegneria del Software.
